\begin{frame}
	\onslide<1->{
	\frametitle{Statistics} 
	\framesubtitle{Random Variables}
}
\begin{itemize}
  \onslide<2->{
\item In measure-theoretic terms, a ra\textbf{ndom variable} 
is simply a real-valued measurable function on the sample space $\Omega$.   
}
  \onslide<3->{
\item A \textbf{measurable function} is one such that the inverse image of any 
open set in $\Real$ is measurable (i.e. is a member of the $\sigma$-algebra).   
}
  \onslide<4->{
\item Some elementary probability texts are not that strict and instead consider 
any real-valued function of the sample space to be a random variable.  
}
\end{itemize}
\end{frame}

\begin{frame}
  \onslide<1->{
	\frametitle{Statistics} 
	\framesubtitle{Random Variables}
}
\begin{itemize}
  \onslide<2->{
\item We use capital letters at the end of the alphabet to denote random variables:  $X,Y,Z, ...$ 
(This helps us distinguish them from ordinary real-valued variables.)  
}
  \onslide<3->{
\item Later we will use lowercase letters so as not to confuse them with certain transforms.  
}
\end{itemize}
\end{frame}
\begin{frame}
  \onslide<1->{
	\frametitle{Statistics} 
	\framesubtitle{Random Variables}
}
\begin{itemize}
  \onslide<2->{
\item Note that random variables are not really variables in any ordinary sense -- they are functions.  Thus, if $\omega\in \Omega$, and $X$ is a random 
variable on $\Omega$, then $X(\omega)\in\Real$, i.e., $X(\omega)$ is a real number.  
}
  \onslide<3->{
\item Also note that $X^{-1}(A)=\set{\omega\in\Omega | X(\omega)\in A}$ is 
always an event in the event algebra by definition.  
}
\end{itemize}
\end{frame}
\begin{frame}
  \onslide<1->{
	\frametitle{Statistics} 
	\framesubtitle{Random Variables}
}
\begin{itemize}
  \onslide<2->{
\item Random variables provide a way to quantify the outcomes of
 experiments  
}
  \onslide<3->{
\item Random variables are important because we are often more
 interested in the consequences of an experiment than in the
 experiment itself  
}
  \onslide<4->{
\item For example, a gambler is usually more interested in his
 wins and losses than in the particular game he's playing  
}
\end{itemize}
\end{frame}
\begin{frame}
  \onslide<1->{
	\frametitle{Statistics} 
	\framesubtitle{Random Variables}
}
\begin{itemize}
  \onslide<2->{
\item Moreover, random variables are not really random  
}
  \onslide<3->{
\item They are deterministic functions from the space of all possible outcomes
 to the real numbers.  If you give them a particular outcome to operate on,
 they return the same real value every time.  
}
  \onslide<4->{
\item What gives them the appearance of randomness is really the fact that you don't 
really know what those values are.  Usually, 
you only know (or assume you know) certain statistical properties like 
distributions (to be discussed soon).  
}
\end{itemize}
\end{frame}
\begin{frame}
  \onslide<1->{
	\frametitle{Statistics} 
	\framesubtitle{Random Variables}
}
\begin{itemize}
  \onslide<2->{
\item The "randomness" of random variables stems from the
 probabilistic nature of the underlying sample space.  Usually, we
 can't say much about the individual outcomes in a sample space.
 We can only measure the size of events. (The connection between the
 "randomness" of the underlying space and the random variables will
 be made by means of probability distributions.)  
}
  \onslide<3->{
\item It is because random variables are functions that we find
 ourselves having to deal with infinite dimensional vector
 spaces, i.e., function spaces.  
}
\end{itemize}
\end{frame}
\begin{frame}
  \onslide<1->{
	\frametitle{Statistics} 
	\framesubtitle{Equality and Equivalence of Random Variables}
}
\begin{itemize}
  \onslide<2->{
\item Since random variables are functions, strict equality of
 random variables is the same as the equality of functions.  Two
 functions are equal if and only if they have equal values when
 evaluated at the same point:  $X=Y\Leftrightarrow
 X(\omega)=Y(\omega)$ for all $\omega\in\Omega$.  
}
  \onslide<3->{
\item However, when you take probability into account, things can
 get a little fuzzier.  We say that two random variables $X$ and
 $Y$ are equal \textit{\textbf{almost surely}} (a.s.) if $P\set{\omega\in\Omega:X(\omega)\neq
 Y(\omega)}=0$.  
}
  \onslide<4->{
\item In other words, $X=Y \text{a.s.}$, if the set of all points at
 which $X$ and $Y$ differ is insignificant, i.e., of measure 0.  
}
\end{itemize}
\end{frame}
\begin{frame}
  \onslide<1->{
	\frametitle{Statistics} 
	\framesubtitle{Simple Random Variables}
}
\begin{itemize}
  \onslide<2->{
\item If a random variable $X$ takes on only finitely many
 values, then $X$ is called a \textit{\textbf{simple random variable}}.
 \item If $X$ is a simple random variable and
 $\set{a_1,\cdots,a_n}$ are its values, let $A_i=\set{\omega:X(\omega)=a_i}$
 for each $i=1,\ldots,n$.\\  
}
  \onslide<3->{
Then $X$ can bewritten as a
 sum:
\begin{equation}\label{simple_rv}
  X=\sum_{i=1}^{n}a_i\chi_{A_i}
\end{equation}
where $\chi_{A_i}$ is the characteristic or indicator function on $A_i$ defined by\\
 $\chi_{A_i}(\omega)=
  \begin{cases}
    1 & \text{if } \omega\in A_i, \\
    0 & \text{otherwise}.
  \end{cases}$
}
\end{itemize}
\end{frame}

\begin{frame}
  \onslide<1->{
	\frametitle{Statistics} 
	\framesubtitle{Example:  Simple Random Variable}
}
\begin{itemize}
  \onslide<2->{
\item Double coin toss:  
}
  \onslide<3->{
  \begin{itemize}
  \item Sample Space: $\Omega=\set{HH,HT,TH,TT}$
}
  \onslide<4->{
\item $\sigma$-algebra: $\mathcal{M}=\set{\emptyset, \set{HH},\set{HT,TH,TT},\Omega}$  
}
  \onslide<5->{
\item Probability Measure: $P(A)=$number of elements in $A$ divided by total number of elements in $\Omega$
\end{itemize}  
}
\end{itemize}
\end{frame}

\begin{frame}
  \onslide<1->{
	\frametitle{Statistics} 
	\framesubtitle{Example:  Simple Random Variable}
}
\begin{itemize}
  \onslide<2->{
\item Let $A=\set{HT,TH,TT}$ and define $X=\chi_A$.  Then
\begin{align*}
X(HH)&=0\\
X(HT)&=1\\
X(TH)&=1\\
X(TT)&=1
\end{align*}  
}
\end{itemize}
\end{frame}

\begin{frame}
  \onslide<1->{
	\frametitle{Statistics} 
	\framesubtitle{Example:  Simple Random Variable}
}
\begin{itemize}
  \onslide<2->{
\item Let $B=\set{HH}$ and define $Y=-3\chi_A + 0.5\chi_B$  
}
  \onslide<3->{
\item $Y$ is a simple random variable.  
}
  \onslide<4->{
\item We compute:  
\begin{align*}
  Y(HH)&=-3\chi_A(HH) + 0.5\chi_B(HH)\\
&=-3\cdot 0 + 0.5\cdot 1\\
&=0.5
\end{align*}
}
\end{itemize}
\end{frame}

\begin{frame}
  \onslide<1->{
	\frametitle{Statistics} 
	\framesubtitle{Random Variables form a Vector Space}
}
\begin{itemize}
  \onslide<2->{
\item Let $RV(\Omega)$ denote the set of all random variables on $\Omega$.
 If $X,Y\in RV(\Omega)$ and $a,b\in \Real$, then $aX+bY\in
 RV(\Omega)$.  
}
  \onslide<3->{
\item Scalar multiplication:  Recall that $(aX+bY)(\omega):=a\cdot X(\omega)+b\cdot
 Y(\omega)$, where $\cdot$ denotes ordinary real multiplication.
 }
   \onslide<4->{
 \item For brevity we will drop the $\cdot$ and merely
 juxtapose.  Likewise, we will drop $\Omega$ from $RV(\Omega)$
 and merely write $RV$, assuming implicitly that the underlying sample space is
 $\Omega$. 
}
  \onslide<5->{
  \item Thus, $RV(\Omega)$ becomes a vector space over the reals.  It is an infinite dimensional vector space.\\
  In fact, after we define a norm and inner products on $RV(\Omega)$ using statistical expectation, 
  we will see that it becomes a Hilbert space (modulo almost surely equivalence).
  }
\end{itemize}
\end{frame}
